\clearpage\pagestyle{otMain}\setlength{\otBlocEspacement}{2.5mm}\fontsize{11}{15}\selectfont\corrigesChapitre{16}{Parallélogrammes et aires}
\begin{blocLivre}{cor-16-1}\begin{corrige}{1}{Côtés opposés}\demarche{Associer les côtés opposés.}
Les côtés opposés ont la même longueur : $\res{CD=7\ \mathrm{cm}}$ et $\res{AD=4\ \mathrm{cm}}$.
\label{fig-cor-16-1}\figureGeo[corrige]{para-1}
\end{corrige}\end{blocLivre}
\begin{blocLivre}{cor-16-2}\begin{corrige}{2}{Diagonales et milieu}\demarche{Utiliser le milieu commun.}
O est le milieu des deux diagonales : $AO=10:2=\res5$ cm et $BO=6:2=\res3$ cm.
\label{fig-cor-16-2}\figureGeo[corrige]{para-2}
\end{corrige}\end{blocLivre}
\begin{blocLivre}{cor-16-3}\begin{corrige}{3}{Un rectangle}\demarche{Utiliser une propriété caractéristique.}
Un parallélogramme possédant un angle droit est un $\res{\text{rectangle}}$. Les autres angles sont aussi droits.
\label{fig-cor-16-3}\figureGeo[corrige]{para-3}
\end{corrige}\end{blocLivre}
\begin{blocLivre}{cor-16-4}\begin{corrige}{4}{Aire directe}\demarche{Associer base et hauteur.}
$A=8\times3=\res{24\ \mathrm{cm}^2}$. On utilise la hauteur perpendiculaire, pas le côté oblique.
\label{fig-cor-16-4}\figureGeo[corrige]{para-4}
\end{corrige}\end{blocLivre}
\begin{blocLivre}{cor-16-5}\begin{corrige}{5}{Un losange}\demarche{Combiner les données avec les propriétés.}
Les côtés opposés sont égaux. Avec deux côtés consécutifs égaux, les quatre le sont : ABCD est un $\res{\text{losange}}$.
\label{fig-cor-16-5}\figureGeo[corrige]{para-5}
\end{corrige}\end{blocLivre}
\begin{blocLivre}{cor-16-6}\begin{corrige}{6}{Construire le quatrième sommet}\demarche{Construire à partir de la définition.}
Tracer par A la parallèle à (BC), et par C la parallèle à (AB). Leur intersection est D. Les deux paires de côtés opposés sont alors parallèles.
\label{fig-cor-16-6}\figureGeo[corrige]{para-6}
\end{corrige}\end{blocLivre}
\begin{blocLivre}{cor-16-7}\begin{corrige}{7}{Construire par les diagonales}\demarche{Utiliser une propriété caractéristique des diagonales.}
Construire O, milieu de [AC], puis D, symétrique de B par rapport à O. Les diagonales [AC] et [BD] ont le même milieu, donc ABCD est un parallélogramme.
\label{fig-cor-16-7}\figureGeo[corrige]{para-7}
\end{corrige}\end{blocLivre}
\begin{blocLivre}{cor-16-8}\begin{corrige}{8}{Rectangle par diagonales}\demarche{Établir les hypothèses avant de conclure.}
Le milieu commun prouve d’abord le parallélogramme. Ses diagonales sont égales : c’est un $\res{\text{rectangle}}$. Rien ne garantit ici un carré.
\label{fig-cor-16-8}\figureGeo[corrige]{para-8}
\end{corrige}\end{blocLivre}
\begin{blocLivre}{cor-16-9}\begin{corrige}{9}{Losange par diagonales}\demarche{Enchaîner deux propriétés.}
Le milieu commun prouve le parallélogramme ; la perpendicularité des diagonales prouve ensuite le $\res{\text{losange}}$.
\label{fig-cor-16-9}\figureGeo[corrige]{para-9}
\end{corrige}\end{blocLivre}
\begin{blocLivre}{cor-16-10}\begin{corrige}{10}{Un carré à reconnaître}\demarche{Croiser les deux caractérisations.}
Diagonales égales : c’est un rectangle. Diagonales perpendiculaires : c’est un losange. Comme les deux propriétés sont réunies, c’est un $\res{\text{carré}}$.
\label{fig-cor-16-10}\figureGeo[corrige]{para-10}
\end{corrige}\end{blocLivre}
\begin{blocLivre}{cor-16-11}\begin{corrige}{11}{Hauteur inconnue}\demarche{Retrouver un facteur du produit.}
$12h=54$, donc $h=54:12=\res{4{,}5\ \mathrm{cm}}$. Cette longueur est perpendiculaire à la base.
\label{fig-cor-16-11}\figureGeo[corrige]{para-11}
\end{corrige}\end{blocLivre}
\begin{blocLivre}{cor-16-12}\begin{corrige}{12}{Unités différentes}\demarche{Convertir les longueurs avant le produit.}
$35$ cm $=0{,}35$ m. L’aire vaut $1{,}2\times0{,}35=\res{0{,}42\ \mathrm{m}^2}$, soit $\res{4200\ \mathrm{cm}^2}$.
\label{fig-cor-16-12}\figureGeo[corrige]{para-12}
\end{corrige}\end{blocLivre}
\begin{blocLivre}{cor-16-13}\begin{corrige}{13}{Une figure composée}\demarche{Additionner des aires de parties disjointes.}
Le rectangle a une aire de 32 cm$^2$. Le triangle a une aire de $8\times3:2=12$ cm$^2$. Le total vaut $\res{44\ \mathrm{cm}^2}$. Les parties ne se recouvrent pas.
\label{fig-cor-16-13}\figureGeo[corrige]{para-13}
\end{corrige}\end{blocLivre}
\begin{blocLivre}{cor-16-14}\begin{corrige}{14}{Une hypothèse manquante}\demarche{Vérifier toutes les conditions de la propriété.}
Les diagonales sont perpendiculaires, mais n’ont pas le même milieu. Les quatre côtés n’ont pas la même longueur. Le quadrilatère n’est pas un losange. Il manque l’hypothèse qu’il est déjà un parallélogramme, ou que les diagonales se coupent en leur milieu.
\label{fig-cor-16-14}\figureGeo[corrige]{para-14}
\end{corrige}\end{blocLivre}
\begin{blocLivre}{cor-16-15}\begin{corrige}{15}{Deux méthodes pour une aire}\demarche{S’appuyer sur les demi-diagonales.}
Chaque triangle rectangle a des côtés perpendiculaires de 5 et 3 cm, donc une aire de $5\times3:2=7{,}5$ cm$^2$. Les quatre donnent $\res{30\ \mathrm{cm}^2}$. On retrouve $(10\times6):2$.
\label{fig-cor-16-15}\figureGeo[corrige]{para-15}
\end{corrige}\end{blocLivre}
